% ============================================================================
%  Rozdział: Badania porównawcze i wnioski
%  Plik przeznaczony do samodzielnej kompilacji lub włączenia przez \input{}
%  do głównego dokumentu pracy magisterskiej.
% ============================================================================
\documentclass[12pt,a4paper]{report}
\usepackage[utf8]{inputenc}
\usepackage[T1]{fontenc}
\usepackage[polish]{babel}
\usepackage{amsmath,amssymb}
\usepackage{booktabs}
\usepackage{multirow}
\usepackage{graphicx}
\usepackage{enumitem}
\usepackage{hyperref}
\usepackage{geometry}
\geometry{margin=2.5cm}

\begin{document}

\chapter{Badania porównawcze metod kalibracji i wnioski}
\label{chap:badania-wnioski}

\section{Wprowadzenie}
Celem przeprowadzonych badań było ilościowe porównanie trzech podejść do wyznaczania
parametrów kamery: dwóch wytrenowanych sieci neuronowych oraz klasycznego, referencyjnego
algorytmu kalibracji opartego na detekcji szachownicy. Badanie miało odpowiedzieć na pytanie,
w jakim stopniu regresja parametrów kamery za pomocą sieci głębokich jest w stanie zastąpić
lub uzupełnić metodę klasyczną, a także wskazać mocne i słabe strony każdego z podejść.

Porównaniu poddano następujące metody:
\begin{itemize}[leftmargin=*]
  \item \textbf{ResNet-50} (\texttt{calibration\_net\_best\_resnet}) --- model jednoklatkowy
        przewidujący wektor $[f/W, k_1, k_2]$;
  \item \textbf{CNN+LSTM} (\texttt{calibration\_net\_best\_lstm}) --- model sekwencyjny (okno
        $5$ klatek) przewidujący pełny wektor dziewięciu parametrów;
  \item \textbf{metoda klasyczna} --- funkcja \texttt{cv2.calibrateCamera} realizująca model
        Browna--Conrady'ego, tożsama z algorytmem użytym w skrypcie \texttt{projectScript.ipynb}.
\end{itemize}

Oba modele sieciowe pochodziły z punktów kontrolnych zapisanych po pierwszej epoce treningu
(strata walidacyjna odpowiednio $0.0082$ dla CNN+LSTM i $0.0133$ dla ResNet-50), co należy
uwzględnić przy interpretacji wyników.

\section{Metodyka}

\subsection{Zbiory testowe}
Eksperyment przeprowadzono na dwóch niezależnych zbiorach danych o rozdzielczości
$640 \times 480$ pikseli:
\begin{enumerate}[leftmargin=*]
  \item \textbf{Zbiór syntetyczny z prawdą podstawową} (\texttt{data/aug\_camera\_test\_seq/}) ---
        $20$ sekwencji po $50$ klatek, z których każda posiada plik \texttt{camera\_params.yaml}
        zawierający rzeczywiste wartości parametrów ($f_x, f_y, c_x, c_y, k_1, k_2, p_1, p_2, k_3$).
        Wewnętrzne narożniki szachownicy: $8 \times 5$.
  \item \textbf{Realne pary stereo} (\texttt{src/calibration\_cv/pairs/}) --- zdjęcia
        rzeczywistej szachownicy podzielone na kamerę \emph{lewą} i \emph{prawą} (po $49$ zdjęć,
        narożniki $9 \times 6$). Zbiór nie posiada prawdy podstawowej, dlatego referencją jest
        wynik kalibracji klasycznej.
\end{enumerate}

\subsection{Ujednolicenie parametrów}
Ze względu na odmienne konwencje normalizacji zastosowane podczas treningu, predykcje sieci
przeliczano do wspólnej jednostki (piksele obrazu oryginalnego). Dla modelu ResNet-50
ogniskową odtwarzano jako $f = f_{\text{norm}} \cdot W$, gdzie $W = 640$; dla modelu CNN+LSTM
parametry denormalizowano przez rozmiar wejścia sieci ($224$ px). Współczynniki dystorsji
porównywano bezpośrednio.

\subsection{Miary jakości}
Do oceny zastosowano trzy miary błędu względem prawdy podstawowej: średni błąd bezwzględny
(MAE), medianę błędu bezwzględnego (MedAE) oraz medianę błędu względnego ogniskowej.
Dla parametru $\theta$ i $N$ sekwencji:
\begin{equation}
  \text{MAE} = \frac{1}{N}\sum_{i=1}^{N} \big|\hat{\theta}_i - \theta_i\big|,
  \qquad
  \text{MedAE} = \operatorname{median}_i \big|\hat{\theta}_i - \theta_i\big|.
\end{equation}
Medianę wprowadzono jako miarę \emph{odporną na wartości odstające}, ponieważ klasyczna
kalibracja bywa niestabilna dla obrazów silnie zniekształconych.

\section{Wyniki --- zbiór z prawdą podstawową}
Tabela~\ref{tab:wyniki-gt} przedstawia zbiorcze błędy estymacji ogniskowej $f_x$ dla
$20$ sekwencji.

\begin{table}[htbp]
  \centering
  \caption{Błąd estymacji ogniskowej $f_x$ na zbiorze syntetycznym z prawdą podstawową.}
  \label{tab:wyniki-gt}
  \begin{tabular}{@{}lccc@{}}
    \toprule
    \textbf{Metoda} & \textbf{MedAE $f_x$ [px]} & \textbf{mediana bł. wzgl.} & \textbf{MAE $f_x$ [px]} \\
    \midrule
    klasyczna  & \textbf{17.6} & \textbf{2.7\,\%}  & 321$^\ast$ \\
    ResNet-50  & 66.8          & 8.7\,\%           & 76.2 \\
    CNN+LSTM   & 496.9         & 70.8\,\%          & 484.7 \\
    \bottomrule
  \end{tabular}

  \vspace{2pt}
  {\footnotesize $^\ast$ Wysoka wartość MAE wynika z dwóch sekwencji odstających
   (\texttt{sequence\_015}, \texttt{sequence\_016}).}
\end{table}

Najważniejsze obserwacje:
\begin{itemize}[leftmargin=*]
  \item \textbf{Metoda klasyczna} osiąga najniższy \emph{medianowy} błąd ogniskowej
        (ok. $2.7\,\%$), co potwierdza jej wysoką dokładność przy pewnej detekcji narożników.
        Znaczna różnica pomiędzy MAE ($321$ px) a MedAE ($17.6$ px) ujawnia jej wrażliwość:
        dla dwóch silnie zniekształconych sekwencji detekcja zawiodła, a estymowana ogniskowa
        osiągnęła wartości $f_x \approx 3300$--$3700$ px (błąd RMS $> 1.9$).
  \item \textbf{ResNet-50} jest najlepszym estymatorem ogniskowej \emph{bez wzorca
        kalibracyjnego} --- z pojedynczego obrazu uzyskuje medianę błędu ok. $8.7\,\%$, mimo
        że był trenowany wyłącznie na dystorsji radialnej, bez perspektywy i składowej
        tangencjalnej.
  \item \textbf{CNN+LSTM} zwraca niemal stałą wartość ogniskowej ($\approx 205$ px, błąd
        $\approx 70\,\%$). Punkt kontrolny zapisano po pierwszej epoce, a model nie nauczył się
        jeszcze regresji i przewiduje wartość zbliżoną do średniej zbioru treningowego.
\end{itemize}

\subsection{Współczynniki dystorsji}
W przypadku współczynników $k_1, k_2$ zaobserwowano istotną prawidłowość: klasyczna
kalibracja odzyskuje dystorsję z \emph{odwrotnym znakiem} względem prawdy podstawowej.
Wynika to z faktu, że generator danych nakłada zniekształcenie jako odwzorowanie
\emph{wprost}, podczas gdy \texttt{cv2.calibrateCamera} modeluje przekształcenie odwrotne.
Dla trudnych sekwencji magnituda $k_2$ bywała skrajnie zawyżona (MedAE $k_2 \approx 7.1$).
Obie sieci przewidywały natomiast $k_1, k_2$ bliskie zeru, co przy niewielkich wartościach
prawdziwych daje niski błąd liczbowy, lecz nie świadczy o poprawnym modelowaniu dystorsji.

\section{Wyniki --- realne pary stereo}
Na zbiorze rzeczywistych zdjęć (brak prawdy podstawowej) klasyczną kalibrację przyjęto za
referencję. Wyniki dla kamery lewej i prawej zestawiono w tabeli~\ref{tab:wyniki-pairs}.

\begin{table}[htbp]
  \centering
  \caption{Estymacja parametrów dla realnych par stereo (referencja: metoda klasyczna).}
  \label{tab:wyniki-pairs}
  \begin{tabular}{@{}llcccc@{}}
    \toprule
    \textbf{Kamera} & \textbf{Metoda} & \textbf{$f_x$ [px]} & \textbf{$f_y$ [px]} & \textbf{$k_1$} & \textbf{$k_2$} \\
    \midrule
    \multirow{3}{*}{lewa}
      & klasyczna & 330.2 & 330.7 & $-0.325$ & $0.112$ \\
      & ResNet-50 & 361.3 & 361.3 & $-0.051$ & $-0.058$ \\
      & CNN+LSTM  & 223.8 & 297.5 & $-0.325$ & $-0.026$ \\
    \midrule
    \multirow{3}{*}{prawa}
      & klasyczna & 325.5 & 326.8 & $-0.331$ & $0.126$ \\
      & ResNet-50 & 375.3 & 375.3 & $-0.036$ & $-0.033$ \\
      & CNN+LSTM  & 223.7 & 297.7 & $-0.324$ & $-0.026$ \\
    \bottomrule
  \end{tabular}
\end{table}

Wnioski z tej części:
\begin{itemize}[leftmargin=*]
  \item \textbf{ResNet-50} najlepiej odwzorowuje ogniskową względem referencji
        ($361$/$375$ px wobec $330$/$325$ px).
  \item \textbf{CNN+LSTM}, mimo słabej estymacji ogniskowej, trafnie odtwarza współczynnik
        $k_1$ ($-0.325$ wobec referencyjnego $\approx -0.33$).
  \item Wyniki dla obu kamer są wzajemnie spójne dla każdej z metod, co potwierdza
        jednorodność zestawu stereo.
\end{itemize}
Wymaga podkreślenia, że metoda klasyczna dawała stosunkowo wysoki błąd RMS ($\approx 2.75$)
na zdjęciach rzeczywistych, co wskazuje na obecność szumu pomiarowego niewystępującego w
danych syntetycznych.

\section{Dyskusja i wnioski końcowe}
Przeprowadzone badania prowadzą do następujących wniosków:
\begin{enumerate}[leftmargin=*]
  \item Do dokładnej kalibracji przy dostępnym wzorcu szachownicy \textbf{najlepsza pozostaje
        metoda klasyczna} --- pod warunkiem pewnej detekcji narożników osiąga medianowy błąd
        ogniskowej rzędu $2.7\,\%$. Jej słabością jest jednak brak odporności: pojedyncze
        nieudane detekcje drastycznie zaniżają jakość, co uzasadnia stosowanie miar medianowych.
  \item \textbf{ResNet-50} jest obiecującym estymatorem ogniskowej z pojedynczego obrazu, bez
        konieczności użycia wzorca kalibracyjnego (błąd $\approx 8.7\,\%$). Czyni to sieć
        użyteczną w zastosowaniach, w których szachownica jest niedostępna.
  \item \textbf{CNN+LSTM} w obecnej postaci (trening ograniczony do jednej epoki) nie jest
        praktycznie użyteczny i wymaga znacznie dłuższego procesu uczenia.
  \item \textbf{Żadna z sieci nie odwzorowuje poprawnie współczynników dystorsji.} Stanowi to
        główny kierunek dalszych prac: dłuższy trening, szerszy i bardziej realistyczny zakres
        parametrów oraz ujednolicenie konwencji znaku dystorsji pomiędzy generatorem danych a
        modelem kalibracji.
\end{enumerate}
Podsumowując, metody uczenia głębokiego oferują komplementarne możliwości względem klasycznej
kalibracji: umożliwiają zgrubną estymację ogniskowej bez wzorca, lecz nie osiągają jeszcze jej
dokładności ani wiarygodności w zakresie zniekształceń optycznych.

\end{document}
